% !TeX root = ../regression_logistique.tex
\chapter{Parcours, prérequis et conventions}
\label{ch:contrat}

\section{Public et prérequis}

Ce cours s'adresse à un étudiant ou un praticien qui veut construire un
classifieur linéaire de bout en bout, formules comprises, plutôt qu'en appeler
une implémentation. Les notions propres au modèle --- cote, logit, entropie
croisée, gradient, hessienne, convexité, un-contre-tous --- sont l'objet du
cours ; le reste est supposé acquis.

\begin{table}[htbp]
\centering
\small
\begin{tabularx}{\textwidth}{@{\hspace{1mm}}%
  >{\raggedright\arraybackslash}m{34mm}%
  >{\raggedright\arraybackslash}X%
  >{\raggedright\arraybackslash}m{35mm}@{\hspace{1mm}}}
{\sffamily\scriptsize\bfseries\color{ink!58}DOMAINE} &
{\sffamily\scriptsize\bfseries\color{ink!58}ATTENDU À L'ENTRÉE} &
{\sffamily\scriptsize\bfseries\color{ink!58}RAPPEL DISPONIBLE} \\
\noalign{\vskip2pt}
\paneltablerow{statistique descriptive}
  {moyenne, médiane, écart-type, valeur manquante}
  {\annexeref{ann:prerequis}}
\noalign{\vskip2pt}
\paneltablerow{analyse}
  {exponentielle, logarithme, fonction croissante}
  {\annexeref{ann:prerequis}}
\noalign{\vskip2pt}
\paneltablerow{géométrie}
  {repère, droite, produit scalaire, norme}
  {\annexeref{ann:prerequis}, \annexeref{ann:geometrie}}
\noalign{\vskip2pt}
\paneltablerow{calcul différentiel}
  {dérivée d'une fonction d'une variable}
  {\annexeref{ann:prerequis}}
\noalign{\vskip2pt}
\paneltablerow[warning]{programmation}
  {Python : fonctions, listes, boucles, imports}
  {aucun}
\end{tabularx}
\caption{Prérequis. Les mathématiques couvrent la partie~I et les annexes de
démonstration ; la programmation intervient dans la partie~II.}
\label{tab:prerequis}
\end{table}

\textbf{Objectifs évaluables de la partie~I.}

\begin{panelenumerate}
  \panelitem{calculer le score et la probabilité d'une observation à partir des
    coefficients et des statistiques de préparation ;}
  \panelitem{écrire la perte d'entropie croisée sous une forme stable en double
    précision, et la contrôler en $z=\pm1000$ ;}
  \panelitem{prédire le signe et l'ordre de grandeur d'une composante du
    gradient ;}
  \panelitem[warning]{distinguer un protocole d'évaluation valide d'un protocole
    entaché de fuite du jeu de test ;}
  \panelitem[warning]{lire une matrice de confusion et l'assortir d'une
    incertitude ;}
  \panelitem[rulecolor]{énoncer la portée exacte de la convexité ;}
  \panelitem[rulecolor]{dire à quelles conditions les sorties d'un classifieur
    un-contre-tous se lisent comme des probabilités de classe.}
\end{panelenumerate}

La partie~II écrit le programme correspondant et suppose la partie~I. Les
annexes portent les démonstrations et supposent la partie~I ainsi que les
rappels de l'\annexeref{ann:prerequis}. Le fil principal énonce les résultats
employés ; les encadrés \emph{approfondissement} renvoient à leurs
démonstrations.

\section{Carte et conventions}

\begin{figure}[htbp]
\centering
\begin{graphique}{Vue d'ensemble de la régression logistique en six panneaux.
La rangée supérieure décrit une prédiction pour des coefficients fixés : les
variables standardisées placent une observation dans un plan séparé par la
droite de score nul ; la sigmoïde transforme ce score en probabilité ; un seuil
produit la classe prédite. La rangée inférieure décrit l'apprentissage et le
contrôle : les deux courbes de perte dépendent de l'étiquette observée ; les
lignes de niveau du coût guident la mise à jour des coefficients dans la
direction opposée au gradient ; un nuage tenu hors de l'apprentissage sert à
comparer les classes prédites aux classes observées.}
\begin{tikzpicture}[
  x=1cm,y=1cm,
  panel/.style={draw=ink!28,fill=white,rounded corners=2pt,line width=.45pt},
  flow/.style={-{Latex[length=1.7mm]},line width=.65pt,ink!58},
  axis/.style={-{Latex[length=1.35mm]},line width=.45pt,ink!60},
  guide/.style={densely dashed,line width=.4pt,ink!38},
  tag/.style={font=\tiny\bfseries,text=accent!88!ink},
  head/.style={font=\scriptsize\bfseries,text=ink!88}]

% Cadres et titres
\path[panel] (0,5.1) rectangle (4.6,9.1);
\path[fill=accentlight,rounded corners=2pt] (0,8.52) rectangle (4.6,9.1);
\node[head,anchor=west] at (.18,8.81) {SCORE ET FRONTIÈRE};
\node[tag,anchor=east] at (4.42,8.81)
  {ch.~\ref{ch:standardisation}--\ref{ch:nuage}};

\path[panel] (5.1,5.1) rectangle (9.7,9.1);
\path[fill=accentlight,rounded corners=2pt] (5.1,8.52) rectangle (9.7,9.1);
\node[head,anchor=west] at (5.28,8.81) {LIEN LOGISTIQUE};
\node[tag,anchor=east] at (9.52,8.81) {ch.~\ref{ch:probabilite}};

\path[panel] (10.2,5.1) rectangle (14.8,9.1);
\path[fill=accentlight,rounded corners=2pt] (10.2,8.52) rectangle (14.8,9.1);
\node[head,anchor=west] at (10.38,8.81) {DÉCISION};
\node[tag,anchor=east] at (14.62,8.81) {ch.~\ref{ch:classifieur}};

\path[panel] (0,0) rectangle (4.6,4.2);
\path[fill=rulelight,rounded corners=2pt] (0,3.62) rectangle (4.6,4.2);
\node[anchor=west,font=\fontsize{6.4}{7}\selectfont\bfseries,text=ink!88]
  at (.18,3.91) {COÛT ET OPTIMISATION};
\node[tag,anchor=east,text=rulecolor!90!ink] at (4.42,3.91)
  {ch.~\ref{ch:descente}--\ref{ch:modele}};

\path[panel] (5.1,0) rectangle (9.7,4.2);
\path[fill=warninglight,rounded corners=2pt] (5.1,3.62) rectangle (9.7,4.2);
\node[head,anchor=west] at (5.28,3.91) {PERTE INDIVIDUELLE};
\node[tag,anchor=east,text=warning!90!ink] at (9.52,3.91)
  {ch.~\ref{ch:erreur}};

\path[panel] (10.2,0) rectangle (14.8,4.2);
\path[fill=gridgray!22,rounded corners=2pt] (10.2,3.62) rectangle (14.8,4.2);
\node[head,anchor=west] at (10.38,3.91) {HORS ÉCHANTILLON};
\node[tag,anchor=east] at (14.62,3.91) {ch.~\ref{ch:predire}};

% Dépendances entre les six représentations
\draw[flow] (4.6,7.05) -- node[above,font=\tiny] {$z_i$} (5.1,7.05);
\draw[flow] (9.7,7.05) -- node[above,font=\tiny] {$p_i$} (10.2,7.05);
\draw[flow] (7.4,5.1) -- node[left,font=\tiny] {$p_i$} (7.4,4.2);
\draw[flow] (5.1,2.05) -- (4.6,2.05);
\draw[flow] (2.3,4.2) -- node[right,font=\tiny] {$w^{(t+1)}$} (2.3,5.1);
\draw[flow] (12.5,5.1) -- node[right,font=\tiny,align=left]
  {modèle\\figé} (12.5,4.2);

% 1. Variables standardisées, score et frontière
\node[font=\tiny,text=ink!68] at (2.3,8.28)
  {$x_j=(a_j-\mu_j)/s_j$ ; statistiques d'apprentissage};
\draw[axis] (.62,5.73) -- (4.18,5.73) node[below,font=\tiny] {$x_1$};
\draw[axis] (.62,5.73) -- (.62,8.12) node[left,font=\tiny] {$x_2$};
\draw[line width=.75pt,ink!78] (1.05,7.95) -- (3.95,5.91);
\node[font=\tiny,fill=white,inner sep=1pt,text=ink!75] at (3.62,6.26)
  {$z=0$};
\foreach \p in {(1.02,6.08),(1.38,6.46),(1.72,6.23),(2.06,6.68),
                (2.34,6.18)}
  \fill[warning] \p circle (1.45pt);
\foreach \p in {(2.18,7.45),(2.62,7.21),(2.96,7.67),(3.35,7.91),
                (3.67,7.48)}
  \node[fill=accent,minimum size=3.1pt,inner sep=0pt] at \p {};
\draw[-{Latex[length=1.6mm]},line width=.8pt,accent]
  (2.42,6.92) -- (3.05,7.81) node[right,font=\tiny] {$w$};
\node[font=\scriptsize] at (2.3,5.36)
  {$z_i=w_0+w_1x_{i1}+w_2x_{i2}$};

% 2. Transformation monotone du score en probabilité
\begin{scope}[shift={(7.4,5.78)},x=.54cm,y=2.18cm]
  \draw[axis] (-3.25,0) -- (3.35,0) node[below,font=\tiny] {$z$};
  \draw[axis] (0,-.04) -- (0,1.12) node[left,font=\tiny] {$p$};
  \draw[guide] (-3.05,.5) -- (3.05,.5);
  \draw[guide] (0,0) -- (0,.5);
  \draw[line width=1.05pt,accent]
    plot[samples=70,domain=-3.1:3.1] (\x,{1/(1+exp(-\x))});
  \fill[accent] (0,.5) circle (1.6pt);
  \node[font=\tiny,anchor=east] at (-.08,.54) {$0{,}5$};
  \node[font=\tiny,anchor=north] at (0,-.04) {$0$};
\end{scope}
\node[font=\scriptsize] at (7.4,5.36)
  {$p_i=\sigma(z_i)=\dfrac{1}{1+e^{-z_i}}$};

% 3. Seuil de décision sur l'échelle des probabilités
\path[fill=warninglight] (10.62,6.25) rectangle (12.5,7.76);
\path[fill=accentlight] (12.5,6.25) rectangle (14.38,7.76);
\draw[line width=.55pt,ink!55] (10.62,7.02) -- (14.38,7.02);
\draw[line width=.7pt,ink!75] (12.5,6.15) -- (12.5,7.9);
\node[font=\tiny,anchor=north] at (10.62,6.98) {$0$};
\node[font=\tiny,anchor=north] at (12.5,6.14) {$\tau$};
\node[font=\tiny,anchor=north] at (14.38,6.98) {$1$};
\node[font=\scriptsize\bfseries,text=warning] at (11.56,7.48)
  {$\widehat y=0$};
\node[font=\scriptsize\bfseries,text=accent] at (13.44,7.48)
  {$\widehat y=1$};
\foreach \x in {10.94,11.31,11.86,12.16}
  \fill[warning] (\x,6.73) circle (1.35pt);
\foreach \x in {12.73,13.08,13.61,14.07}
  \node[fill=accent,minimum size=3pt,inner sep=0pt] at (\x,6.73) {};
\node[font=\scriptsize] at (12.5,5.58)
  {$\widehat y_i=\mathbf 1[p_i\geq\tau]$};

% 4. Les deux pertes possibles pour une même valeur du score
\begin{scope}[shift={(7.4,.68)},x=.5cm,y=.72cm]
  \draw[axis] (-3.25,0) -- (3.35,0) node[below,font=\tiny] {$z$};
  \draw[axis] (-3.08,0) -- (-3.08,3.35)
    node[left,font=\tiny] {$\ell$};
  \draw[line width=.9pt,accent]
    plot[samples=70,domain=-3:3] (\x,{ln(1+exp(-\x))});
  \draw[line width=.9pt,warning]
    plot[samples=70,domain=-3:3] (\x,{ln(1+exp(\x))});
  \node[font=\tiny,text=accent,anchor=west] at (-2.85,2.75) {$y=1$};
  \node[font=\tiny,text=warning,anchor=east] at (2.85,2.75) {$y=0$};
\end{scope}
\node[font=\tiny,text=warning!85!ink] at (7.4,3.34)
  {$y_i$ est observé sur le jeu d'apprentissage};
\node[font=\scriptsize] at (7.4,.25)
  {$\ell(z_i,y_i)=\operatorname{softplus}(z_i)-y_i z_i$};

% 5. Surface du coût, gradient et mise à jour des coefficients
\begin{scope}[rotate around={-18:(2.24,2.03)}]
  \draw[line width=.45pt,rulecolor!42] (2.24,2.03) ellipse (1.62cm and 1.02cm);
  \draw[line width=.55pt,rulecolor!58] (2.24,2.03) ellipse (1.12cm and .7cm);
  \draw[line width=.65pt,rulecolor!76] (2.24,2.03) ellipse (.62cm and .38cm);
\end{scope}
\fill[rulecolor] (2.24,2.03) circle (1.7pt);
\node[font=\tiny,text=rulecolor!80!ink,anchor=west] at (2.36,1.93)
  {$w^\star$};
\fill[ink!78] (3.66,3.12) circle (1.6pt);
\fill[ink!58] (3.02,2.65) circle (1.35pt);
\draw[-{Latex[length=1.7mm]},line width=.9pt,rulecolor]
  (3.66,3.12) -- (3.02,2.65)
  node[midway,above left,font=\tiny,text=rulecolor!90!ink] {$-\nabla J$};
\draw[-{Latex[length=1.5mm]},line width=.7pt,rulecolor!75]
  (3.02,2.65) -- (2.57,2.34);
\node[font=\tiny,text=ink!62] at (1.25,3.25)
  {$J(w)=\dfrac{1}{n}\sum_i\ell_i(w)$};
\node[font=\scriptsize] at (2.3,.38)
  {$w^{(t+1)}=w^{(t)}-\alpha\nabla J\!\left(w^{(t)}\right)$};

% 6. Contrôle sur des observations absentes de l'ajustement
\draw[axis] (10.62,.72) -- (14.4,.72) node[below,font=\tiny] {$x_1$};
\draw[axis] (10.62,.72) -- (10.62,3.37) node[left,font=\tiny] {$x_2$};
\draw[line width=.75pt,ink!72] (10.98,3.25) -- (14.18,.93);
\foreach \p in {(10.98,1.08),(11.38,1.42),(11.76,1.21),(12.15,1.72)}
  \fill[warning] \p circle (1.45pt);
\foreach \p in {(12.45,2.45),(12.91,2.2),(13.36,2.73),(13.84,2.47)}
  \node[fill=accent,minimum size=3.1pt,inner sep=0pt] at \p {};
\draw[warning,line width=.65pt] (12.49,2.41) -- (12.41,2.49)
  (12.41,2.41) -- (12.49,2.49);
\draw[accent,line width=.65pt] (12.19,1.68) -- (12.11,1.76)
  (12.11,1.68) -- (12.19,1.76);
\node[font=\tiny,align=center,text=ink!68] at (12.5,.27)
  {comparer $\widehat y_i$ à $y_i$ : matrice de confusion,\\
   exactitude, rappel, précision et incertitude};
\end{tikzpicture}
\end{graphique}
\caption{Vue d'ensemble du modèle. La rangée supérieure transforme une
observation en classe prédite pour des coefficients fixés. La rangée inférieure
relie l'étiquette observée à la perte, à l'optimisation des coefficients et au
contrôle sur un jeu tenu hors de l'ajustement.}
\label{fig:carte}
\end{figure}
\FloatBarrier

Le trajet supérieur de la figure~\ref{fig:carte} est celui d'une prédiction. Le
trajet inférieur n'intervient que pendant l'apprentissage ; l'évaluation
réutilise ensuite les statistiques de préparation et les coefficients figés,
sans les réajuster sur le jeu de test.

\textbf{Notation.} Les symboles sont déclarés dans le tableau de
l'\annexeref{ann:notation}, avec le glossaire français--anglais. Deux collisions
y sont signalées : $\sigma$ désigne l'écart-type d'une colonne comme la fonction
logistique, selon le contexte ; la cote standard usuellement notée $z$ est notée
$x$ ici, la lettre $z$ étant réservée au score.

\textbf{États du modèle.} Les coefficients dépendent du point d'arrêt de la
descente. Quatre états portent un nom, et chaque tableau ou figure indique
lequel il emploie.

\begin{table}[htbp]
\centering
\begin{tabular}{cl>{\raggedright\arraybackslash}p{62mm}}
\toprule
nom & coefficients & obtenu par \\
\midrule
\Mzero  & $(0\,;0\,;0)$ & point de départ de toute descente \\
\Mquatre & $(-4{,}745\,;-3{,}454\,;+3{,}469)$
  & $400$ tours à $\alpha=1$ sur les $1\,600$ observations ; figures de la
    partie~I \\
\Mstop & $(-5{,}176\,;-3{,}746\,;+3{,}787)$
  & même descente poursuivie jusqu'au critère de stagnation, $983$ tours ;
    partie~II \\
\Movr & quatre jeux de trois nombres
  & quatre descentes un-contre-tous, une par classe \\
\bottomrule
\end{tabular}
\caption{États du modèle. \Mquatre et \Mstop diffèrent par la norme des
coefficients : leurs rapports coïncident à trois millièmes près
(ch.~\ref{ch:descente}).}
\label{tab:etats}
\end{table}

\textbf{Provenance des chiffres.} Statistiques descriptives, coefficients, coûts
et taux proviennent du fichier \texttt{dataset\_train.csv} du sujet DSLR de 42 et
des scripts déposés dans \texttt{scripts/} ; les traces employées par les figures
sont dans \texttt{data/}, et \texttt{scripts/verifie\_chiffres.py} les recalcule
toutes. Le jeu de données est synthétique et les personnages qu'il nomme sont
fictifs.

\textbf{Vocabulaire.} Une perte, un coût ou une erreur qualifient une
\emph{prédiction} portant sur une observation : le document écrit « perte
associée à cette prédiction ». La règle maintient la responsabilité du côté du
modèle lorsque les données décrivent des personnes. Le raisonnement général
emploie \emph{observation}, \emph{variable} et \emph{classe} ; \emph{élève},
\emph{note} et les maisons sont réservés aux exemples tirés du fichier.

\textbf{Limites du développement.} Régularisation, sélection de variables,
validation croisée, calibration des probabilités, régression multinomiale
complète, optimisation de second ordre, coûts asymétriques et dérive de
distribution sont situés dans le cycle d'utilisation, sans démonstration ni
implémentation complète. Leur rôle et la limite qu'ils imposent au résultat
sont indiqués aux points concernés.
